\documentclass[11pt]{article}
\usepackage[a4paper,margin=1in]{geometry}
\usepackage[T1]{fontenc}
\usepackage{lmodern,microtype}
\usepackage{amsmath,amssymb,amsthm,mathtools}
\usepackage{booktabs,enumitem}
\usepackage{xcolor}
\usepackage[numbers,sort&compress]{natbib}
\usepackage[colorlinks=true,linkcolor=blue!55!black,citecolor=blue!55!black]{hyperref}

\newtheorem{theorem}{Theorem}[section]
\newtheorem{proposition}[theorem]{Proposition}
\theoremstyle{remark}
\newtheorem{remark}[theorem]{Remark}

\title{A Cloud-Extracted Trace-Moment Atlas through Order Twelve}
\author{Liang Wang}
\date{July 2026}

\begin{document}
\maketitle

\begin{abstract}
We compute the first twelve power traces of the 32 fine and Haar-coarse
operators underlying the RH-222 cloud atlas.  From each trace we subtract the
Perron mode, parity mode, and shell-complete selected cloud.  This produces
384 cloud-extracted trace moments without constructing the ill-conditioned
Riesz projectors of RH-232.

For the logarithmic unit-disk seminorm
\[
 J_{12}(\tau)=\sum_{n=2}^{12}\frac{|\tau_n|}{n},
\]
the batch maximum is $0.07593$.  On the fine range $\sigma\le0.005$ it is
$0.01067$.  Across all endpoints and orders $2$--$12$, the largest observed
root rate $|\tau_n|^{1/n}$ is $0.35989$.  At the finest left endpoint the
orders two through six have moduli $2.07\times10^{-2}$,
$9.67\times10^{-4}$, $1.30\times10^{-5}$, $8.97\times10^{-6}$, and
$5.47\times10^{-8}$.

The atlas supplies strong finite-order evidence for a trace-based relative
determinant route.  It does not control order thirteen or the infinite tail,
so Gate A remains open.
\end{abstract}

\section{Trace extraction without a complement matrix}

For a finite matrix $A$ with listed peripheral values $p,q$ and selected
cloud $C$, define
\begin{equation}\label{eq:tau}
 \tau_n(A;C)=\operatorname{tr}(A^n)-p^n-q^n
 -\sum_{\lambda\in C}\lambda^n.
\end{equation}
This quantity depends only on the finite matrix trace and selected eigenvalue
multiset.  It is therefore insensitive to the norm of the associated
spectral projector.

\begin{proposition}[Reality under conjugacy closure]\label{prop:reality}
If $A$ is real, $p,q\in\mathbb R$, and $C$ is closed under complex
conjugation with multiplicity, then $\tau_n(A;C)\in\mathbb R$ for every
$n\ge1$.
\end{proposition}

\begin{proof}
The trace of a real matrix power is real.  Every nonreal pair
$\lambda,\overline\lambda$ contributes
$\lambda^n+\overline\lambda^n=2\operatorname{Re}(\lambda^n)$, and real
cloud values contribute real powers.
\end{proof}

The archived imaginary parts are consequently roundoff diagnostics rather
than independent spectral data.  Their scale stays near ordinary floating
arithmetic, providing a useful check that shell completion and multiplicity
bookkeeping were preserved.

\begin{proposition}[Finite determinant jet]\label{prop:jet}
Let $R_C(z)$ be the projection-free complementary product from RH-234.  Near
the origin,
\[
 \log R_C(z)=-\sum_{n=2}^{m}\frac{\tau_n(A;C)}{n}z^n+O(z^{m+1}).
\]
Consequently
\[
 \sup_{|z|\le R}
 \left|\sum_{n=2}^{m}\frac{\tau_n(A;C)}{n}z^n\right|
 \le \sum_{n=2}^{m}\frac{|\tau_n(A;C)|}{n}R^n.
\]
\end{proposition}

\begin{proof}
Expand each regularized eigenvalue factor using
$\log((1-w)e^w)=-\sum_{n\ge2}w^n/n$, then subtract the selected cloud powers.
The inequality is the triangle inequality.
\end{proof}

\section{Sparse computation}

The fine matrices have dimensions from $128$ to $4096$; the coarse matrices
have half those dimensions.  We form sparse powers recursively and take the
diagonal sum.  No unresolved eigenvalue decomposition is used.  At the finest
left endpoint, the twelfth power has approximately $1.49\times10^7$ stored
entries, still below the dense $4096^2$ count.

The numerical traces are exact for the frozen double-precision sparse
matrices up to ordinary floating arithmetic.  They are not interval
enclosures for a continuum operator.

\section{Moment decay across orders}

\begin{table}[ht]
\centering
\scriptsize
\begin{tabular}{rrrrr}
\toprule
$n$ & max all $|\tau_n|$ & max fine $|\tau_n|$
& all rate & fine rate\\
\midrule
2  & $1.2952\times10^{-1}$ & $2.0670\times10^{-2}$ & $0.3599$ & $0.1438$\\
3  & $2.9813\times10^{-2}$ & $1.0758\times10^{-3}$ & $0.3101$ & $0.1025$\\
4  & $3.6811\times10^{-3}$ & $1.6452\times10^{-4}$ & $0.2463$ & $0.1133$\\
5  & $7.6942\times10^{-4}$ & $8.9640\times10^{-6}$ & $0.2384$ & $0.0978$\\
6  & $6.1019\times10^{-4}$ & $2.9136\times10^{-6}$ & $0.2912$ & $0.1195$\\
7  & $3.0091\times10^{-4}$ & $4.4089\times10^{-7}$ & $0.3140$ & $0.1236$\\
8  & $1.1350\times10^{-4}$ & $2.2973\times10^{-8}$ & $0.3213$ & $0.1110$\\
9  & $3.7095\times10^{-5}$ & $2.8206\times10^{-9}$ & $0.3219$ & $0.1122$\\
10 & $1.0776\times10^{-5}$ & $1.9058\times10^{-10}$ & $0.3186$ & $0.1067$\\
11 & $2.7665\times10^{-6}$ & $1.9500\times10^{-11}$ & $0.3124$ & $0.1063$\\
12 & $5.9728\times10^{-7}$ & $1.9573\times10^{-12}$ & $0.3029$ & $0.1058$\\
\bottomrule
\end{tabular}
\caption{Complete determinant-relevant order atlas.  ``Fine'' means
$\sigma\le0.005$, and each rate is the corresponding maximum modulus raised
to the power $1/n$.}
\end{table}

The sequence is not monotonically decreasing in every endpoint or every
order.  Complex phase cancellation and rank changes produce local spikes.
The relevant observation is that the entire finite envelope stays far below
the failed Hilbert--Schmidt budget of RH-229 \citep{WangRH229,WangRH235}.

Three summaries answer different questions.  The root rate
$|\tau_n|^{1/n}$ tests compatibility with geometric decay; the jet norm
$J_{12}$ bounds a truncated logarithm on the unit disk; and the raw modulus
records the actual coefficient scale.  A small jet norm can be dominated by
order two even if later root rates rise, which is why all three are archived
rather than a single fitted exponent.

\section{Scale behavior}

The fine-scale jet norms are small but not monotone.  For example, the left
channel rises again at $\sigma=0.00125$ because the fixed rank schedule changes
by another shell.  This prevents a direct Cauchy conclusion from the current
selection.  It motivates two separate tests:
\begin{enumerate}[leftmargin=2em]
\item compare the same finite jet between the fine and coarse discretizations;
\item choose the shell prefix by a trace tolerance rather than a predetermined
rank.
\end{enumerate}
RH-237 and RH-238 perform these tests.

\section{Numerical protocol and reproducibility boundary}

The 384 cases are $32$ endpoints times orders one through twelve.  Only the
$352$ cases at orders two through twelve enter the regularized logarithm.
The order-one values are retained as a bookkeeping check on peripheral
subtraction.  For each endpoint the recursion starts with the frozen sparse
matrix and repeatedly multiplies by that same matrix; the trace is read from
the diagonal after every multiplication.  The cloud contribution is then
subtracted from the stored shell-complete eigenvalue list.

This procedure avoids the two unstable operations diagnosed earlier in the
route: no left/right overlap is inverted, and no full unresolved spectrum is
diagonalized.  Its cost is memory growth in the sparse powers.  At dimension
$4096$, the order-twelve power contains about $1.49\times10^7$ stored
entries, which explains why twelve is a computational horizon rather than a
mathematically distinguished order.

The archive verifies conjugacy closure, trace-array lengths, the unit-disk
majorants, and consistency with the RH-235 second moments.  It does not bound
sparse multiplication roundoff or discretization error relative to a
continuum transfer operator.

\section{The missing theorem}

An order-$12$ jet can neither exclude large coefficients at later orders nor
guarantee a zero-free disk for the full complementary determinant.  The
necessary strengthening is a uniform estimate
\[
 |\tau_n(\sigma)|\le Mq^n,
 \qquad n\ge2,
\]
with $q<1/R$ on a target disk $|z|\le R$.  RH-240 proves that this all-order
condition is sufficient.  The present paper supplies only its first eleven
nontrivial data points.

At least three tails remain compatible with the table: continued geometric
decay, a delayed high-order spike, or a small tail produced by extracting too
many regular modes into the moving cloud.  The first would advance the
determinant route, the second would block the proposed unit disk, and the
third would give the wrong normalization despite excellent numerical bounds.
This trichotomy is why an all-order envelope and a coefficient anchor are
separate theorem obligations.

Gates A--E remain open, and no arithmetic spectral identification is claimed.

\bibliographystyle{plainnat}
\bibliography{references}
\end{document}
